\documentclass[0pt]{article}
\usepackage[margin=0.5in]{geometry}
\usepackage{listings}
\usepackage{pgfplots}
\usepackage{xcolor}
\lstset{
	language = C,
	frame=tb, % draw a frame at the top and bottom of the code block
	tabsize=4, % tab space width
	showstringspaces=false, % don't mark spaces in strings
	commentstyle=\color{green}, % comment color
	keywordstyle=\color{blue},  % keyword color
	stringstyle=\color{red}  % string color
}
\begin{document}
	\begin{center}
		\begin{LARGE}
			\textcolor{red}{\underline{SORTING ALGORITHMS}}
		\end{LARGE}
		\\
		\leavevmode \\
	\end{center}
	% Bubble Sort Starts here 
	\begin{flushleft}
		\textbf{\textcolor{red}{1. Bubble Sort}}
		\leavevmode \\
		\begin{description}
			\item[Approach - ] Bubble sort is based on the idea of repeatedly comparing pairs of adjacent elements and then swapping their positions if they exist in the wrong order. This is the most time consuming sorting Algorithms as the time complexity of the Bubble sort goes like this :  
			\leavevmode \\
			\begin{description}
				\item[Worst Case] - \(O(n^2)\)
				\item[Average Case] - \(O(n^2)\)
				\item[Best Case] - \(O(n)\)
			\end{description}
			\leavevmode \\
			\item[Program - ] \leavevmode \\
			\leavevmode \\
			\begin{lstlisting}
			void bubble_sort(int *a,int n)
			{
				for (int i=0; i<n;i++)
				{
					for(int j=0; j<n-i-1;i++)
					{
						if(a[j]>a[j+1])
						{
							swap(a[i],a[j]);
						}
					}
				}
			}
			\end{lstlisting}
			\item[Plot - ] \leavevmode \\
		\end{description}
	\end{flushleft}
	\begin{center}
		\begin{tikzpicture}
		\begin{axis}[
		axis lines=middle,
		xmin=0, xmax=11000,
		ymin=0, ymax=0.6,
		xtick=\empty, ytick=\empty
		]
		\addplot [only marks] table {
			0 0
			1000 0.009556
			2000 0.02371
			3000 0.04547
			4000 0.074976
			5000 0.113558
			6000 0.165204
			7000 0.236573
			8000 0.305851
			9000 0.399738
			10000 0.49556
		};
		\end{axis}
		\end{tikzpicture}
	\end{center}
	\leavevmode \\
	% Bubble Sort Ends here
	
	
	
	
	
	
	
	
	
	
	
	
	
	
	
	
	
	
	
	
	
	
	
	
	
	
	
	
	
	
	
	
	
	
	
	
	
	% Selection Sort Starts here 
	\begin{flushleft}
		\textbf{\textcolor{red}{2. Selection Sort}}
		\leavevmode \\
		\begin{description}
			\item[Approach - ] The Selection sort algorithm is based on the idea of finding the minimum or maximum element in an unsorted array and then putting it in its correct position in a sorted array. Time complexity of the Bubble sort goes like this :  
			\leavevmode \\
			\begin{description}
				\item[Worst Case] - \(O(n^2)\)
				\item[Average Case] - \(O(n^2)\)
				\item[Best Case] - \(O(n^2)\)
			\end{description}
			\leavevmode \\
			\item[Program - ] \leavevmode \\
			\leavevmode \\
			\begin{lstlisting}
			void selection_sort(int *a,int n)
			{
				for(int i=0;i<n-1;i++)
				{
					imin=i;
					for(int j=i+1;j<n;j++)
					{
						if(a[j]<a[imin]) 
							imin=j;
					}
					swap(a[imin],a[i])
				}
			}
			\end{lstlisting}
			\item[Plot - ] \leavevmode \\
		\end{description}
	\end{flushleft}
	\begin{center}
		\begin{tikzpicture}
		\begin{axis}[
		axis lines=middle,
		xmin=0, xmax=11000,
		ymin=0, ymax=0.6,
		xtick=\empty, ytick=\empty
		]
		\addplot [only marks] table {
			0 0
			1000 0.003479
			2000 0.013566
			3000 0.027378
			4000 0.041152
			5000 0.062114
			6000 0.085499
			7000 0.119053
			8000 0.143657
			9000 0.180674
			10000 0.22365
		};
		\end{axis}
		\end{tikzpicture}
	\end{center}
	\leavevmode \\
	% Selection Sort Ends here
	
	
	
	
	
	
	
	
	
	
	
	
	
	
	
	
	
	
	
	
	
	
	
	
	
	
	
	
	% Insertion Sort Starts here 
	\begin{flushleft}
		\textbf{\textcolor{red}{3. Insertion Sort}}
		\leavevmode \\
		\begin{description}
			\item[Approach - ] Insertion sort is based on the idea that one element from the input elements is consumed in each iteration to find its correct position i.e, the position to which it belongs in a sorted array. Time complexity of the Bubble sort goes like this :  
			\leavevmode \\
			\begin{description}
				\item[Worst Case] - \(O(n^2)\)
				\item[Average Case] - \(O(n^2)\)
				\item[Best Case] - \(O(n)\)
			\end{description}
			\leavevmode \\
			\item[Program - ] \leavevmode \\
			\leavevmode \\
			\begin{lstlisting}
			void insertion_sort(int *a,int n)
			{
				for(int i=1;i<n;i++){
					value = a[i];
					hole =i;
					while(hole>0&&a[hole-1]>value){
						a[hole] = a[hole-1];
						hole--;}
					a[hole] =value;
				}
			}
			\end{lstlisting}
			\item[Plot - ] \leavevmode \\
		\end{description}
	\end{flushleft}
	\begin{center}
		\begin{tikzpicture}
		\begin{axis}[
		axis lines=middle,
		xmin=0, xmax=11000,
		ymin=0, ymax=0.6,
		xtick=\empty, ytick=\empty
		]
		\addplot [only marks] table {
			0 0
			1000 0.003592
			2000 0.009168
			3000 0.019161
			4000 0.029052
			5000 0.039449
			6000 0.055332
			7000 0.068427
			8000 0.089401
			9000 0.113026
			10000 0.135115
		};
		\end{axis}
		\end{tikzpicture}
	\end{center}
	\leavevmode \\
	% Insertion Sort Ends here
	
	
	
	
	
	
	
	
	
	
	
	
	
	
	
	
	
	
	
	
	
	
	
	
	
	
	
	
	
	
	% Merge Sort Starts here 
	\begin{flushleft}
		\textbf{\textcolor{red}{4. Merge Sort}}
		\leavevmode \\
		\begin{description}
			\item[Approach - ] Merge sort is a divide-and-conquer algorithm based on the idea of breaking down a list into several sub-lists until each sublist consists of a single element and merging those sublists in a manner that results into a sorted list. Time complexity of the Bubble sort goes like this :  
			\leavevmode \\
			\begin{description}
				\item[Worst Case] - \(O(nlogn)\)
				\item[Average Case] - \(O(nlogn)\)
				\item[Best Case] - \(O(nlogn)\)
			\end{description}
			\leavevmode \\
			\item[Program - ] \leavevmode \\
			\leavevmode \\
			\begin{lstlisting}
			void merge_sort(int *a,int n)
			{
				int mid,*l,*r;
				if(size==1) { return 1;}
				mid = size/2;
				l = new int [mid];
				r = new int [size-mid];
				for(int i=0;i<mid;i++) l[i] = a[i];
				for(int i=mid;i<size;i++) r[i-mid] = a[i];
				MergeSort(l,mid);
				MergeSort(r,size-mid);
				Merge(l,r,a,mid,size-mid);
			}
			\end{lstlisting}
			\item[Plot - ] \leavevmode \\
		\end{description}
	\end{flushleft}
	\begin{center}
		\begin{tikzpicture}
		\begin{axis}[
		axis lines=middle,
		xmin=0, xmax=11000,
		ymin=0, ymax=0.01,
		xtick=\empty, ytick=\empty
		]
		\addplot [only marks] table {
			0 0
			1000 0.000853
			2000 0.001374
			3000 0.00196
			4000 0.002169
			5000 0.003228
			6000 0.005415
			7000 0.007569
			8000 0.00666
			9000 0.0061199
			10000 0.008269
		};
		\end{axis}
		\end{tikzpicture}
	\end{center}
	\leavevmode \\
	% Merge Sort Ends here
	
	
	
	
	
	
	
	
	
	
	
	
	
	
	
	
	
	
	
	
	
	
	% Quick Sort Starts here 
	\begin{flushleft}
		\textbf{\textcolor{red}{5. Quick Sort}}
		\leavevmode \\
		\begin{description}
			\item[Approach - ] Quick sort is based on the divide-and-conquer approach based on the idea of choosing one element as a pivot element and partitioning the array around it such that: Left side of pivot contains all the elements that are less than the pivot element Right side contains all elements greater than the pivot Time complexity of the Bubble sort goes like this :  
			\leavevmode \\
			\begin{description}
				\item[Worst Case] - \(O(n^2)\)
				\item[Average Case] - \(O(nlogn)\)
				\item[Best Case] - \(O(nlogn)\)
			\end{description}
			\leavevmode \\
			\item[Program - ] \leavevmode \\
			\leavevmode \\
			\begin{lstlisting}
			void quick_sort(int *a,int start,int end)
			{
				int pIndex;
				if(start<end)
				{
					pIndex=partition(a,start,end);
					QuickSort(a,start,pIndex-1);
					QuickSort(a,pIndex+1,end);
				}
			}
			\end{lstlisting}
			\item[Plot - ] \leavevmode \\
		\end{description}
	\end{flushleft}
	\begin{center}
		\begin{tikzpicture}
		\begin{axis}[
		axis lines=middle,
		xmin=0, xmax=110000,
		ymin=0, ymax=0.01,
		xtick=\empty, ytick=\empty
		]
		\addplot [only marks] table {
			0 0
			10000 0.003608
			20000 0.006167
			30000 0.006628
			40000 0.008968
			50000 0.007334
			60000 0.006411
			70000 0.005522
			80000 0.004264
			90000 0.005855
			100000 0.00607
		};
		\end{axis}
		\end{tikzpicture}
	\end{center}
	\leavevmode \\
	% Quick Sort Ends here
	
	
	
	
	
	
	
	
	
	
	
	
	
	
	
	
	
	
	
	
	
	
	
	
	
	% Counting Sort Starts here 
	\begin{flushleft}
		\textbf{\textcolor{red}{6. Counting Sort}}
		\leavevmode \\
		\begin{description}
			\item[Approach - ] In Counting sort, the frequencies of distinct elements of the array to be sorted is counted and stored in an auxiliary array, by mapping its value as an index of the auxiliary array. Time complexity of the Bubble sort goes like this :  
			\leavevmode \\
			\begin{description}
				\item[Worst Case] - \(O(n+k)\)
				\item[Average Case] - \(O(n+k)\)
				\item[Best Case] - \(O(n+k)\)
			\end{description}
			\leavevmode \\
			\item[Program - ] \leavevmode \\
			\leavevmode \\
			\begin{lstlisting}
			void counting_sort(int *arr,int n,int RANGE)
			{
				int count[RANGE]={0};
				int out[n];
				for(i=0;i<n;i++) ++count[arr[i]];
				for(i=1;i<RANGE;i++) count[i]+=count[i-1];
				for(i=n-1;i>=0;i--){
				out[count[arr[i]]-1]=arr[i];
				--count[arr[i]];
				}
				for(i=0;i<n;i++) arr[i]=out[i];
			}
			\end{lstlisting}
			\item[Plot - ] \leavevmode \\
		\end{description}
	\end{flushleft}
	\begin{center}
		\begin{tikzpicture}
		\begin{axis}[
		axis lines=middle,
		xmin=0, xmax=110000,
		ymin=0, ymax=0.05,
		xtick=\empty, ytick=\empty
		]
		\addplot [only marks] table {
			0 0
			10000 0.001974
			20000 0.016839
			30000 0.004715
			40000 0.020774
			50000 0.009915
			60000 0.02424
			70000 0.015947
			80000 0.032722
			90000 0.023292
			100000 0.040942
		};
		\end{axis}
		\end{tikzpicture}
	\end{center}
	\leavevmode \\
	% Counting Sort Ends here
	
	
	
	
	
	
	
	
	
	
	
	
	
	
	
	
	
	
	
	
	
	
	
	
	
	
	
	
	
	
	
	
	% Heap Sort Starts here 
	\begin{flushleft}
		\textbf{\textcolor{red}{7. Heap Sort}}
		\leavevmode \\
		\begin{description}
			\item[Approach - ] Heaps can be used in sorting an array. In max-heaps, maximum element will always be at the root. Heap Sort uses this property of heap to sort the array.
			\leavevmode \\
			Consider an array \(Arr\) which is to be sorted using Heap Sort.
			
			1) Initially build a max heap of elements in Arr.
			\leavevmode \\
			2) The root element, that is \(Arr[1]\), will contain maximum element of \(Arr\). After that, swap this element with the last element of \(Arr\) and heapify the max heap excluding the last element which is already in its correct position and then decrease the length of heap by one.
			\leavevmode \\
			3) Repeat the step 2, until all the elements are in their correct position.
			Time complexity of the Bubble sort goes like this :  
			\leavevmode \\
			\begin{description}
				\item[Worst Case] - \(O(nlogn)\)
				\item[Average Case] - \(O(nlogn)\)
				\item[Best Case] - \(O(nlogn)\)
			\end{description}
			\leavevmode \\
			\item[Program - ] \leavevmode \\
			\leavevmode \\
			\begin{lstlisting}
			void heap_sort(int *a,int n)
			{
				for (int i = n / 2 - 1; i >= 0; i--)
				heapify(arr, n, i);
				for (int i=n-1; i>=0; i--)
				{
					swap(arr[0], arr[i]);
					heapify(arr, i, 0);
				}	
			}
			\end{lstlisting}
			\item[Plot - ] \leavevmode \\
		\end{description}
	\end{flushleft}
	\begin{center}
		\begin{tikzpicture}
		\begin{axis}[
		axis lines=middle,
		xmin=0, xmax=11000,
		ymin=0, ymax=0.03,
		xtick=\empty, ytick=\empty
		]
		\addplot [only marks] table {
			0 0
			1000 0.00427
			2000 0.001148
			3000 0.006341
			4000 0.004332
			5000 0.010829
			6000 0.010202
			7000 0.016218
			8000 0.016115
			9000 0.022145
			10000 0.019955
		};
		\end{axis}
		\end{tikzpicture}
	\end{center}
	\leavevmode \\
	% Heap Sort Ends here
	
	
oindent\rule{\textwidth}{0.4pt}
\end{document}